\documentclass[11pt,a4paper]{article}
\usepackage[T1]{fontenc}
\usepackage[utf8]{inputenc}
\usepackage[slovak,english]{babel}
\usepackage{geometry,xcolor,array,booktabs,longtable,tabularx,tikz,fancyhdr,titlesec,enumitem,hyperref,lastpage,microtype}
\usepackage[sfdefault]{noto}
\usepackage[most]{tcolorbox}
\usepackage{pdflscape}
\usetikzlibrary{arrows.meta,calc,positioning,shapes.geometric,babel}

\geometry{left=18mm,right=18mm,top=21mm,bottom=20mm,headheight=20pt}
\setlength{\parindent}{0pt}
\setlength{\parskip}{5pt}
\setlength{\emergencystretch}{2em}
\setlist{nosep,leftmargin=*}
\definecolor{limblue}{HTML}{102A43}
\definecolor{limcyan}{HTML}{00A6A6}
\definecolor{limaccent}{HTML}{FF6B4A}
\definecolor{limmint}{HTML}{DFF7F5}
\definecolor{limlight}{HTML}{F1F5F8}
\definecolor{limpaper}{HTML}{FAFCFD}
\definecolor{limgray}{HTML}{627D98}
\hypersetup{colorlinks=true,linkcolor=limcyan,urlcolor=limcyan,pdftitle={LIM M2 --- Architektonická špecifikácia procesora}}
\pagestyle{fancy}
\fancyhf{}
\fancyhead[L]{\colorbox{limblue}{\textcolor{white}{\strut\hspace{2mm}\textbf{LIM M2}\hspace{2mm}}}}
\fancyhead[R]{\small\textcolor{limgray}{PROCESSOR REFERENCE MANUAL / REVISION 0.2}}
\fancyfoot[L]{\small Architektonická príručka LIM M2 --- Revízia 0.2 (Október 2026)}
\fancyfoot[R]{\small \thepage/\pageref{LastPage}}
\renewcommand{\headrulewidth}{0pt}
\renewcommand{\footrulewidth}{0pt}
\titleformat{\section}{\Large\bfseries\color{limblue}}
  {\colorbox{limcyan}{\textcolor{white}{\strut\thesection}}}{0.75em}{}
\titleformat{\subsection}{\large\bfseries\color{limblue}}{\thesubsection}{0.7em}{}
\titlespacing*{\section}{0pt}{16pt}{8pt}

\newcommand{\TBD}{\textcolor{limcyan}{\textbf{TBD}}}
\newcommand{\reg}[1]{\texttt{#1}}
\newcommand{\bits}[1]{\texttt{#1}}
\newcommand{\field}[1]{\textsf{#1}}
\newcommand{\opcodebox}[1]{\tikz[baseline=(n.base)]\node[fill=limblue,text=white,rounded corners=2pt,inner xsep=7pt,inner ysep=3pt,font=\ttfamily\bfseries] (n) {#1};}
\newtcolorbox{limnote}[1]{
  enhanced,breakable,colback=limlight,colframe=limcyan,boxrule=0pt,
  borderline west={3pt}{0pt}{limcyan},arc=1mm,left=4mm,right=4mm,
  top=3mm,bottom=3mm,title={#1},fonttitle=\bfseries\color{limblue},
  coltitle=limblue,attach title to upper={\par\smallskip}}
\newtcolorbox{specbox}{
  enhanced,colback=limpaper,colframe=limlight,boxrule=0.8pt,arc=2mm,
  left=3mm,right=3mm,top=2mm,bottom=2mm}

\newcommand{\InstructionPage}[8]{%
  \clearpage
  \phantomsection\addcontentsline{toc}{subsection}{#1 --- #4}
  \begin{tcolorbox}[enhanced,colback=limblue,colframe=limblue,arc=3mm,
    left=5mm,right=5mm,top=4mm,bottom=4mm]
  \begin{minipage}[t]{0.68\textwidth}
    {\fontsize{30}{34}\selectfont\bfseries\color{white}#1}\par\vspace{1mm}
    {\large\color{white}#4}
  \end{minipage}\hfill
  \begin{minipage}[t]{0.25\textwidth}\raggedleft
    \tikz[baseline=(n.base)]\node[fill=limcyan,text=white,rounded corners=2pt,
      inner xsep=8pt,inner ysep=4pt,font=\ttfamily\bfseries] (n) {#2};
    \par\smallskip{\small\color{white}#3}
  \end{minipage}
  \end{tcolorbox}
  \vspace{3mm}
  \begin{specbox}
  \begin{tabularx}{\dimexpr\textwidth-8mm\relax}{>{\bfseries\color{limblue}}p{31mm}X}
    \toprule
    Mnemonic & \texttt{#1} \\
    Opcode & \bits{#2} (Horných 5 bitov inštrukčného slova) \\
    Formát inštrukcie & 5-bitový opcode, 3-bitové pole \field{cfg} a dynamické argumenty (pri špecializovaných príkazoch argumenty priamo za opkódom) \\
    Trieda operácie & #3 \\
    Stav ratifikácie & Základná špecifikácia schválená; rozvrhnutie bitových polí argumentov vo fáze finalizácie \\
    \bottomrule
  \end{tabularx}
  \end{specbox}
  \subsection*{Účel inštrukcie} #4.
  \subsection*{Architektonická sémantika}
  \begin{tcolorbox}[colback=limmint,colframe=limmint,arc=2mm,boxrule=0pt,left=4mm]
    \ttfamily #5
  \end{tcolorbox}
  \subsection*{Konfigurácia a operandy} #6
  \subsection*{Vplyv na príznakový register (FREG)} #7
  \subsection*{Mikrooperácie dátovej cesty}
  \begin{enumerate}
    \item Multiplexorový výber zdrojových operandov z registrového poľa alebo priamej zbernice.
    \item Odoslanie operácie do vyhradenej funkčnej jednotky (ALU / pamäť / riadenie).
    \item Smerovanie výsledku cez selektor spätného zápisu (MUX 2) do cieľového registra.
    \item Synchrónna aktualizácia stavových príznakov \reg{FREG} a registrových ukazovateľov (\reg{PTREG}, \reg{SREG}).
  \end{enumerate}
  Presné cyklické časovanie a riadiace impulzy: \TBD.
  \subsection*{Poznámky k implementácii} #8
  \vfill{\footnotesize\color{limgray}Polia s označením TBD špecifikujú konfigurovateľné hardvérové voľby.}
}

\begin{document}
\begin{titlepage}
  \pagecolor{limpaper}\color{limblue}
  \begin{tikzpicture}[remember picture,overlay]
    \fill[limblue] (current page.north west) rectangle
      ([xshift=58mm]current page.south west);
    \fill[limcyan] ([xshift=58mm]current page.north west) rectangle
      ([xshift=63mm]current page.south west);
    \draw[limlight,line width=0.7pt] ([xshift=77mm,yshift=-37mm]current page.north west)
      -- ([xshift=-16mm,yshift=-37mm]current page.north east);
  \end{tikzpicture}
  \vspace*{18mm}
  \hspace*{-7mm}\begin{minipage}{50mm}\centering
    {\color{white}\bfseries\large HIGH-PERFORMANCE}\par\smallskip
    {\color{limcyan}\bfseries\small 16-BIT ARCHITECTURE}
  \end{minipage}\par
  \vspace{31mm}
  \hspace*{68mm}\begin{minipage}{100mm}\raggedright
    {\fontsize{38}{42}\selectfont\bfseries LIM\par M2}\par
    \vspace{7mm}
    {\Large\bfseries\color{limcyan}Architektonická príručka procesora}\par
    \vspace{9mm}
    {\large Vysoko deterministická 16-bitová mikroprogramová architektúra s čistou multiplexorovou dátovou cestou a rozšíreným 24-bitovým adresovaním.}\par
    \vspace{18mm}
    \begin{tikzpicture}[node distance=5mm]
      \foreach \x/\n in {0/R0,1.25/R1,2.5/R2,3.75/R3,5/R4,6.25/R5,7.5/R6,8.75/R7}
        \node[draw=limblue,rounded corners=1.5pt,minimum width=10mm,minimum height=8mm,
          fill=white,font=\bfseries\scriptsize] (r\n) at (\x,0) {\n};
      \draw[limcyan,line width=2pt,-{Latex[length=2mm]}] (0,-0.8) -- (8.75,-0.8);
      \draw[limblue,line width=1.2pt,-{Latex[length=2mm]}] (8.75,-1.25) -- (0,-1.25);
      \node[anchor=west,font=\small\bfseries,text=limgray] at (0,-1.85)
        {POINT-TO-POINT DATA ROUTING \enspace·\enspace ZERO BUS CONTENTION};
    \end{tikzpicture}
  \end{minipage}
  \vfill
  \hspace*{68mm}\begin{minipage}{100mm}\raggedright
    {\bfseries\color{limaccent}OFFICIAL REFERENCE SPECIFICATION}\par\smallskip
    {\large Revízia 0.2 --- Október 2026}\par\medskip
    {\small\color{limgray}Špecifikácia jadra LIM M2. Všetky práva vyhradené.}
  \end{minipage}
  \vspace*{10mm}
\end{titlepage}
\nopagecolor

\tableofcontents
\clearpage

\section{Rozsah dokumentu a architektonická filozofia}
Tento dokument predstavuje oficiálnu technickú špecifikáciu a architektonický štandard mikroprocesora \textbf{LIM M2}. 
Architektúra bola navrhnutá s dôrazom na maximálnu spoľahlivosť, matematickú exaktnosť a vysokú taktovaciu frekvenciu.

Základným kameňom dátovej cesty LIM M2 je \textbf{úplné vylúčenie vnútorných zberníc tretieho stavu (Tri-State / Hi-Z)}. Nahradením zdieľaných kapacitných liniek aktívne budenými stromami multiplexorov (Point-to-Point MUX Routing) sa úplne zamedzuje vzniku skratových prúdov, zbernicových konfliktov (bus contention) a parazitných kapacitných oneskorení.

\begin{limnote}{Číslovanie inštrukčného súboru}
Inštrukčná sada LIM M2 využíva súvislý 5-bitový priestor primárnych kódov: 
\texttt{NOP}=0, \texttt{ADD}=1, \ldots, \texttt{MSP}=26. Kódy 27--31 sú vyhradené pre hardvérové rozšírenia novej generácie.
\end{limnote}

\subsection{Dohodnuté označenia}
\begin{tabularx}{\textwidth}{>{\ttfamily}p{28mm}X}
\toprule
Označenie & Architektonický význam \\
\midrule
dst & Cieľový register pre zápis výsledku operácie \\
src, srcA, srcB & Zdrojové registre vstupných operandov \\
imm & Priama skalárna konštanta v tele inštrukcie \\
addr & Efektívna adresa pamäte (16-bitová základná alebo 24-bitová rozšírená) \\
cfg & 3-bitové konfiguračné pole modifikátora operácie \\
ZF, CF, OF & Hardvérové príznaky nuly, prenosu a pretečenia v registri \reg{FREG} \\
\bottomrule
\end{tabularx}

\subsection{Jednotná filozofia formátu inštrukcií}
Architektúra procesora LIM M2 striktne uplatňuje univerzálnu paradigmu: \textbf{najprv sa vždy uvádza cieľový príjemca (kam), a až potom vstupné operandy (odkiaľ)}. To zaručuje čistú, logickú a prehľadnú syntax:
\begin{itemize}
  \item \textbf{Aritmetické operácie (3 operandy):} \texttt{ADD Rd, Rs1, Rs2} znamená $\text{Rd} \leftarrow \text{Rs1} + \text{Rs2}$ (napríklad zápis \texttt{ADD R0, R1, R2} vypočíta $R0 = R1 + R2$); rovnako pre \texttt{SUB}, \texttt{MUL}, \texttt{DIV}.
  \item \textbf{Logické operácie (2 operandy):} \texttt{AND Rd, Rs} znamená $\text{Rd} \leftarrow \text{Rd} \ \text{AND}\ \text{Rs}$ (napríklad \texttt{AND R0, R1} vypočíta $R0 = R0 \ \& \ R1$); rovnako pre \texttt{OR}, \texttt{NAND}, \texttt{NOR}, \texttt{XOR}, \texttt{XNOR}.
  \item \textbf{Bitová inverzia NOT (1 operand, 1-bajtová inštrukcia):} \texttt{NOT Rd} znamená $\text{Rd} \leftarrow \sim\text{Rd}$ (napríklad \texttt{NOT R0} vypočíta $R0 = \sim R0$). Kóduje sa v 1 bajte (\texttt{[7:3]} opkód, \texttt{[2:0]} register Rd), druhý bajt z pamäte sa nenačítava.
  \item \textbf{Medziregistrové presuny:} \texttt{LRR Rd, Rs} vykonáva priame kopírovanie z bodu B (zdroj \reg{Rs}) do bodu A (príjemca \reg{Rd}): $\text{Rd} \leftarrow \text{Rs}$ (napríklad \texttt{LRR R0, R1} vykoná $R0 = R1$).
  \item \textbf{Načítanie z pamäte (LMR):} prvý argument \reg{Rd} určuje cieľový register procesora, do ktorého sa uložia načítané dáta.
\end{itemize}

\section{Architektonický prehľad procesora LIM M2}
\textbf{LIM M2} je všestranný, vysoko výkonný 16-bitový procesor s mikroprogramovým riadením, optimalizovaný na celočíselné výpočty. Medzi kľúčové architektonické prednosti patria:

\begin{itemize}
  \item \textbf{Čisté multiplexorové smerovanie dát:} Vnútorná dátová cesta je kompletne zbavená tretieho stavu. Prenos dát je riadený aktívnymi budičmi cez kaskádu selektorov \textbf{MUX 1}, \textbf{DEST-DMUX} a \textbf{MUX 2}. Architektúra je optimálne prispôsobená moderným technológiám ASIC a rýchlym FPGA.
  \item \textbf{Dvojitý akumulátorový uzol (Dual Accumulator Staging):} Registre \reg{ACCA} a \reg{ACCB} stabilne fixujú operandy pred ALU, čím izolujú prechodové javy výpočtu od hlavného registrového poľa.
  \item \textbf{Hardvérové reťazenie dlhej aritmetiky (ALU Auto-Chaining):} Blok ALU integruje vnútorný zvyškový register (\reg{ALU\_RESIDUAL}). V prípade prenosu ($+1$ pri súčte $>2^{16}-1$), výpožičky ($-1$ pri rozdiele $<0$) alebo vyšších bitov násobenia/delenia sa zvyšok zachytáva hardvérovo a automaticky pripočíta či odpočíta v nasledujúcej matematickej operácii. Tým sa zabezpečuje priebežná viacbajtová aritmetika (32, 48, 64+ bitov) bez samostatných inštrukcií ADC/SBB. Vyprázdnenie reťazca sa vykonáva inštrukciou \texttt{NOP}.
  \item \textbf{Priamy premosťovací kanál (Reg-to-Reg Bypass):} Obchádzková vetva z MUX 1 priamo do MUX 2 umožňuje bleskové kopírovanie medzi registrami bez blokovania aritmetickej jednotky.
  \item \textbf{Lineárne 24-bitové adresovanie pamäte (až 16 MB):} Špecializované 24-bitové registre \reg{PTREG} (ukazovateľ údajov s autoincrementom) a \reg{SREG} (zásobníkový ukazovateľ) umožňujú lineárny prístup k rozsiahlym dátovým blokom a zásobníku až do 16~MB bez segmentačných bariér.
  \item \textbf{Vysoký frekvenčný potenciál:} Odstránenie kapacitného zaťaženia distribuovaných zberníc umožňuje jadru dosahovať základnú taktovaciu frekvenciu \textbf{200~MHz}, pričom v optimalizovanej kremíkovej implementácii škáluje na \textbf{700~MHz a viac}.
\end{itemize}

\clearpage
\begin{landscape}
\thispagestyle{empty}
\subsection{Konceptuálna schéma dátovej cesty (Datapath)}
\begin{center}
\begin{tikzpicture}[
    x=1cm, y=1cm,
    >=Stealth,
    font=\sffamily\scriptsize,
    regnode/.style={
        draw=black!85, rectangle, rounded corners=3pt, minimum width=1.1cm, minimum height=0.75cm,
        line width=1.1pt, fill=white, font=\bfseries\scriptsize, align=center
    },
    block/.style={
        draw=black!85, rectangle, rounded corners=3pt, line width=1.1pt, fill=white,
        font=\bfseries\scriptsize, align=center
    },
    ctrlbox/.style={
        draw=red!80!black, rectangle, rounded corners=7pt, line width=1.4pt,
        fill=red!5, align=center
    },
    membox/.style={
        draw=cyan!75!black, rectangle, rounded corners=6pt, line width=1.3pt,
        fill=cyan!6, align=center
    },
    databus/.style={->, line width=1.3pt, draw=black!90},
    databus_line/.style={line width=1.3pt, draw=black!90},
    ctrlbus/.style={->, line width=1.0pt, draw=red!80!black, dashed},
    ctrlbus_line/.style={line width=1.0pt, draw=red!80!black, dashed},
    bidir/.style={<->, line width=1.2pt, draw=teal!80!black}
]

    \foreach \i/\x in {0/0.6, 1/2.0, 2/3.4, 3/4.8, 4/6.2, 5/7.6, 6/9.0, 7/10.4} {
        \node[regnode] (R\i) at (\x, 12.0) {\textbf{R\i}};
        \coordinate (R\i_in) at (\x - 0.25, 12.38);
        \coordinate (R\i_we) at (\x + 0.25, 12.38);
        \coordinate (R\i_out) at (\x, 11.62);
    }

    \draw[fill=gray!12, draw=black!85, line width=1.2pt]
        (-0.2, 10.6) -- (11.2, 10.6) -- (10, 9.7) -- (1.1, 9.7) -- cycle;
    \node[align=center] at (5.5, 10.15) {\textbf{\normalsize MUX 1}};

    \foreach \i in {0,...,7} {
        \draw[databus] (R\i_out) -- (R\i_out |- 0, 10.6);
    }

    \draw[fill=gray!12, draw=black!85, line width=1.1pt]
        (1.3, 6.9) -- (3.1, 6.9) -- (3.6, 5.9) -- (0.8, 5.9) -- cycle;
    \node[align=center] at (2.2, 6.4) {\textbf{\scriptsize DEST-DMUX}};

    \node[regnode, fill=blue!6, draw=blue!75!black, minimum width=1.2cm, minimum height=0.8cm]
        (ACCA) at (1.3, 4.9) {\textbf{ACCA}};
    \node[regnode, fill=blue!6, draw=blue!75!black, minimum width=1.2cm, minimum height=0.8cm]
        (ACCB) at (3.1, 4.9) {\textbf{ACCB}};

    \draw[databus] (1.3, 5.9) -- (ACCA.north);
    \draw[databus] (3.1, 5.9) -- (ACCB.north);
    \node[font=\tiny\bfseries, fill=white, inner sep=1pt] at (1.05, 5.55) {[0]};
    \node[font=\tiny\bfseries, fill=white, inner sep=1pt] at (3.35, 5.55) {[1]};

    \draw[fill=blue!8, draw=blue!85!black, line width=1.3pt]
        (0.5, 3.6) -- (1.8, 3.6) -- (2.2, 3.2) -- (2.6, 3.6) -- (3.9, 3.6) -- (3.2, 2.0) -- (1.2, 2.0) -- cycle;
    \node[align=center, text=blue!90!black] at (2.2, 2.7) {\textbf{\large ALU}};

    \draw[databus] (ACCA.south) -- (1.3, 3.6);
    \draw[databus] (ACCB.south) -- (3.1, 3.6);
    \node[font=\tiny\bfseries, left] at (1.3, 4.1) {A [15:0]};
    \node[font=\tiny\bfseries, right] at (3.1, 4.1) {B [15:0]};

    \draw[fill=yellow!18, draw=orange!80!black, line width=1.3pt]
        (5.6, 5.8) -- (7.2, 5.8) -- (7.8, 4.4) -- (5.0, 4.4) -- cycle;
    \node[align=center, text=orange!90!black] at (6.4, 5.25) {\textbf{\normalsize MUX 2}};

    \node[font=\tiny\bfseries, fill=white, draw=orange!80!black, rounded corners=1pt, inner sep=1.2pt] at (5.6, 4.68) {ALU};
    \node[font=\tiny\bfseries, fill=white, draw=orange!80!black, rounded corners=1pt, inner sep=1.2pt] at (6.4, 4.68) {REG};
    \node[font=\tiny\bfseries, fill=white, draw=orange!80!black, rounded corners=1pt, inner sep=1.2pt] at (7.2, 4.68) {MEM};

    \node[block, minimum width=3.0cm, minimum height=0.8cm, fill=teal!6, draw=teal!80!black]
        (FLAGS) at (6.4, 1.2) {\textbf{FREG [15:0]}\\[1pt]\scriptsize ZF, CF, OF, EM, \ldots};

    \node[ctrlbox, minimum width=5.4cm, minimum height=5.0cm]
        (CTRL) at (14.1, 6.7) {};

    \node[font=\bfseries\normalsize, text=red!90!black] at (14.1, 8.7) {CONTROL UNIT};

    \node[draw=red!40!black, fill=white, rounded corners=2pt, font=\tiny\bfseries, inner sep=2pt, right]
        at (11.5, 8.3) {IR [15:0]};
    \node[draw=red!70!black, fill=white, rounded corners=2pt, font=\tiny\bfseries, inner sep=2pt, text=red!80!black, right]
        at (11.5, 7.3) {CTRL ALU/ACC};
    \node[draw=red!70!black, fill=white, rounded corners=2pt, font=\tiny\bfseries, inner sep=2pt, text=red!80!black, right]
        at (11.5, 6.2) {SEL\_WB [1:0]};
    \node[draw=teal!70!black, fill=white, rounded corners=2pt, font=\tiny\bfseries, inner sep=2pt, text=teal!90!black, right]
        at (11.5, 5.2) {FLAGS};
    \node[draw=orange!80!black, fill=white, rounded corners=2pt, font=\tiny\bfseries, inner sep=2pt, text=orange!90!black, right]
        at (11.5, 4.5) {MEM/IMM};

    \node[membox, minimum width=5.4cm, minimum height=2.6cm]
        (MEM) at (14.1, 2.1) {\textbf{\normalsize MEMORY}\\[3pt]\tiny RAM / ROM / DEVICES / ELSE};

    \draw[databus, draw=blue!80!black] (12.6, 4.2) -- (12.6, 3.4) node[midway, fill=white, inner sep=1pt, font=\tiny\bfseries, text=blue!80!black] {ADDR};
    \draw[databus] (14.1, 4.2) -- (14.1, 3.4) node[midway, fill=white, inner sep=1pt, font=\tiny\bfseries] {WR\_DATA};
    \draw[bidir, draw=red!80!black] (15.6, 3.4) -- (15.6, 4.2) node[midway, fill=white, inner sep=1pt, font=\tiny\bfseries, text=red!80!black] {RD / WR};

    \draw[databus_line] (5.5, 9.7) -- (5.5, 8.3);
    \filldraw (5.5, 8.3) circle (2pt);

    \draw[databus] (5.5, 8.3) -| (2.2, 6.9) node[pos=0.25, above, font=\tiny\bfseries] {[15:0]};

    \draw[databus] (5.5, 8.3) -- (11.4, 8.3) node[pos=0.45, above, font=\tiny\bfseries] {REG\_DATA [15:0]};

    \filldraw (4.6, 8.3) circle (1.8pt);
    \draw[databus_line] (4.6, 8.3) -- (4.6, 3.6);
    \draw[databus_line] (4.6, 3.6) -- (5.4, 3.6);
    \draw[databus_line] (5.4, 3.6) arc[start angle=180, end angle=0, radius=0.2] -- (6.4, 3.6);
    \draw[databus] (6.4, 3.6) -- (6.4, 4.4);
    \node[font=\tiny\bfseries, fill=white, inner sep=1.2pt, below] at (5.0, 3.55) {[15:0]};

    \draw[databus] (3.4, 2.5) -- (5.6, 2.5) -- (5.6, 4.4);
    \node[font=\tiny\bfseries, fill=white, inner sep=1pt] at (4.5, 2.75) {ALU\_RES [15:0]};

    \draw[databus_line] (11.4, 4.5) -- (10.3, 4.5) -- (10.3, 3.6) -- (9.8, 3.6);
    \draw[databus_line] (9.8, 3.6) arc[start angle=0, end angle=180, radius=0.2] -- (7.2, 3.6);
    \draw[databus] (7.2, 3.6) -- (7.2, 4.4);
    \node[font=\tiny\bfseries, fill=white, inner sep=1pt] at (10.8, 4.7) {MEM/IMM [15:0]};

    \draw[->, line width=1.1pt, draw=teal!80!black] (2.2, 2.0) |- (FLAGS.west);
    \node[font=\tiny\bfseries, fill=white, inner sep=1pt, text=teal!90!black] at (3.5, 1.2) {FLAGS};

    \draw[bidir] (7.9, 1.2) -- (9.6, 1.2) -- (9.6, 5.2) -- (11.4, 5.2);
    \node[font=\tiny\bfseries, fill=white, inner sep=1pt, text=teal!90!black] at (10.4, 5.45) {FLAGS};

    \draw[databus_line] (6.4, 5.8) -- (6.4, 9.0) -- (5.7, 9.0)
                        arc[start angle=0, end angle=180, radius=0.2]
                        -- (-0.8, 9.0) -- (-0.8, 13.5) -- (10.9, 13.5);
    \node[font=\tiny\bfseries, fill=black!85, text=white, inner sep=2.5pt, rounded corners=2pt]
        at (2.2, 9.0) {WB\_DATA [15:0]};

    \foreach \i in {0,...,7} {
        \draw[databus] (R\i_in |- 0, 13.5) -- (R\i_in);
        \filldraw (R\i_in |- 0, 13.5) circle (1.8pt);
    }

    \draw[ctrlbus] (12.8, 9.2) |- (10.6 , 10.15);
    \node[font=\tiny\bfseries, text=red!80!black, fill=white, inner sep=1pt] at (12.8, 9.7) {SEL\_REG [2:0]};

    \draw[ctrlbus] (11.4, 6.2) -| (8.5, 5.3) -- (7.4, 5.3);
    \node[font=\tiny\bfseries, text=red!80!black, fill=white, inner sep=1pt] at (9.8, 6.2) {SEL\_WB [1:0]};

    \draw[ctrlbus_line] (11.4, 7.3) -- (6.6, 7.3)
                        arc[start angle=0, end angle=180, radius=0.2]
                        -- (4.8, 7.3)
                        arc[start angle=0, end angle=180, radius=0.2]
                        -- (4.1, 7.3);

    \draw[ctrlbus_line] (4.1, 7.3) -- (4.1, 4.9);
    \filldraw[red!80!black] (4.1, 7.3) circle (1.4pt);

    \draw[ctrlbus] (4.1, 6.4) -- (3.35, 6.4);
    \filldraw[red!80!black] (4.1, 6.4) circle (1.4pt);
    \node[font=\tiny\bfseries, text=red!80!black, fill=white, inner sep=1pt, above] at (4.1, 6.45) {SEL\_ACC};

    \draw[ctrlbus] (4.1, 4.9) -- (ACCB.east);
    \node[font=\tiny\bfseries, text=red!80!black, fill=white, inner sep=1pt, above] at (4.15, 5.1) {LD\_B};

    \draw[ctrlbus_line] (4.1, 7.3) -- (2.4, 7.3)
                        arc[start angle=0, end angle=180, radius=0.2]
                        -- (0.3, 7.3);

    \draw[ctrlbus] (0.2, 7.3) -- (0.2, 4.9) -- (ACCA.west);
    \node[font=\tiny\bfseries, text=red!80!black, fill=white, inner sep=1pt, left] at (0.4, 5.2) {LD\_A};

    \draw[ctrlbus] (0.2, 4.9) -- (0.2, 2.8) -- (0.8, 2.8);
    \node[font=\tiny\bfseries, text=red!80!black, fill=white, inner sep=1pt, left] at (0.5, 3.3) {ALU\_OP [3:0]};

    \draw[ctrlbus_line] (15.2, 9.2) -- (15.2, 12.8) -- (-0.2, 12.8);
    \node[font=\tiny\bfseries, text=red!80!black, fill=white, inner sep=1pt] at (13.0, 12.8) {REG\_WE [7:0]};
    \foreach \i in {0,...,7} {
        \draw[ctrlbus] (R\i_we |- 0, 12.8) -- (R\i_we);
        \filldraw[red!80!black] (R\i_we |- 0, 12.8) circle (1.4pt);
    }

\end{tikzpicture}
\end{center}
\end{landscape}
\clearpage

\section{Programovo viditeľné registre a stav procesora}
Architektonický stav procesora LIM M2 zahŕňa ortogonálnu sadu funkčných registrov:

\begin{longtable}{>{\ttfamily\bfseries}p{25mm}p{28mm}p{93mm}}
\toprule
Register & Šírka & Architektonický účel a vlastnosti \\
\midrule\endhead
R0--R7 & 16 bitov & Registre všeobecného určenia (GPR). Selektor MUX 1 zabezpečuje bleskové čítanie; zápis prebieha synchronizovane cez zbernicu WB\_DATA pod kontrolou REG\_WE. \\
SREG & 24 bitov & Systémový zásobníkový register (Stack Register). Hardvérová podpora volaní podprogramov a lokálnych rámcov s 24-bitovým adresovaním zásobníka (až 16~MB). \\
PTREG & 24 bitov & Presný ukazovateľový register (Pointer Register). Zabezpečuje lineárne adresovanie pamäte až do 16~MB s voliteľným hardvérovým autoincrementom. \\
FREG & 16 bitov & Stavový a riadiaci register (Status \& Control Register). Konsoliduje aritmetické príznaky ALU a konfiguračné bity jadra. \\
\bottomrule
\end{longtable}

\subsection{Topológia príznakového a riadiaceho registra (FREG)}
Register \reg{FREG} koordinuje činnosť výpočtových a systémových uzlov procesora:

\begin{tabularx}{\textwidth}{>{\bfseries}p{14mm}>{\ttfamily}p{23mm}X}
\toprule
Bit & Mnemonic & Funkčný popis \\
\midrule
0    & ZF   & Zero Flag: nastavený pri nulovom výsledku operácie. \\
1    & CF   & Carry Flag: príznak aritmetického prenosu alebo výpožičky. \\
2    & OF   & Overflow Flag: indikátor znamienkového pretečenia v doplnkovom kóde. \\
3    & EM   & Error Math: indikátor matematických výnimiek (delenie nulou). \\
4--5 & CMP  & Compare Result: stav hardvérového komparátora (00: reset/neporovnané/chyba, 01: menšie, 10: rovné, 11: väčšie). \\
6--7 & CLK  & Clock Mode: voľba deliča systémovej frekvencie jadra. \\
8    & PIE  & Program Interrupt Enable: globálne povolenie softvérových prerušení. \\
9    & MIE  & Master Interrupt Enable: povolenie externých hardvérových prerušení. \\
10   & DWR  & Disable Wait for Ready: nútený zbernicový cyklus bez čakania na potvrdzovací signál. \\
11   & M24  & 24-bit Mode Enable: aktivácia rozšíreného 24-bitového lineárneho adresovania. \\
12   & AIP  & Auto-Increment PTREG: povolenie hardvérového autoincrementu ukazovateľa \reg{PTREG}. \\
13   & DBIT & Displaced Bit: záchytný register pre bit vysunutý pri posunoch. \\
14   & SMOD & Shift Mode: voľba cyklického posunu cez DBIT (0: nulovanie, 1: priechod cez DBIT). \\
15   & INTR & In Interrupt: príznak vykonávania prerušenia (pri 1 sú ďalšie prerušenia zablokované). \\
\bottomrule
\end{tabularx}

\section{Formát strojových inštrukcií}
Inštrukčný súbor LIM M2 je navrhnutý pre rýchle a deterministické dekódovanie. 
Horných 5 bitov je vyhradených pre primárny operačný kód:

\begin{center}
\begin{tikzpicture}[x=0.9cm,y=0.72cm]
  \draw[thick,fill=limblue] (0,0) rectangle (5,1.25);
  \draw[thick,fill=limcyan] (5,0) rectangle (8,1.25);
  \draw[thick,fill=limlight] (8,0) rectangle (17,1.25);
  \node[text=white,font=\bfseries,align=center] at (2.5,0.62) {OPCODE\\5 bitov};
  \node[text=white,font=\bfseries,align=center] at (6.5,0.62) {CFG\\3 bity};
  \node[font=\bfseries,align=center] at (12.5,0.62) {Operandy a argumenty\\dynamická dĺžka};
  \node[anchor=south west,font=\small] at (0,1.32) {$W-1$};
  \node[anchor=south] at (5,1.32) {$W-6$};
  \node[anchor=south] at (8,1.32) {$W-9$};
  \node[anchor=south east,font=\small] at (17,1.32) {0};
\end{tikzpicture}
\end{center}

Pole \field{cfg} (3 bity) selektuje až 8 ortogonálnych režimov inštrukcie, čím poskytuje vysokú flexibilitu bez nežiaduceho rozširovania počtu opkódov.

\section{Prehľadná tabuľka inštrukcií jadra LIM M2}
\small
\begin{longtable}{r>{\ttfamily}p{13mm}>{\ttfamily}p{19mm}p{28mm}p{75mm}}
\toprule
Č. & Kód & Mnemonic & Trieda & Funkčný účel \\
\midrule\endhead
0  & 00000 & NOP  & Riadenie   & Prázdna operácia a vyprázdnenie zvyškového registra ALU \\
1  & 00001 & ADD  & Aritmetika & Súčet s automatickým reťazením prenosu cez zvyšok ALU \\
2  & 00010 & SUB  & Aritmetika & Odčítanie s automatickým reťazením výpožičky cez zvyšok ALU \\
3  & 00011 & MUL  & Aritmetika & Násobenie so zachytením vyšších bitov vo zvyšku ALU \\
4  & 00100 & DIV  & Aritmetika & Delenie s uložením zvyšku do vnútorného registra ALU \\
5  & 00101 & LMR  & Pamäť      & Načítanie z pamäte/vstupu do registra (konfig zdroja 0--3, veľkosť 8/16 bit) \\
6  & 00110 & LRM  & Pamäť      & Zápis dát z registra do pamäte/výstupu (konfig cieľa 0--3, veľkosť 8/16 bit) \\
7  & 00111 & LRR  & Presun     & Priame kopírovanie medzi registrami cez obchádzkový MUX \\
8  & 01000 & CMP  & Porovnanie & Hardvérové porovnanie operandov bez zmeny dát \\
9  & 01001 & JMP  & Skok       & Nepodmienený skok (priamy, nepriamy, 24-bitový) \\
10 & 01010 & JCR  & Skok       & Podmienený skok pri nastavenom príznaku prenosu CF \\
11 & 01011 & JNZ  & Skok       & Podmienený skok pri nenulovom výsledku (ZF = 0) \\
12 & 01100 & JZ   & Skok       & Podmienený skok pri nulovom výsledku (ZF = 1) \\
13 & 01101 & AND  & Logika     & Bitový logický súčin (AND) \\
14 & 01110 & OR   & Logika     & Bitový logický súčet (OR) \\
15 & 01111 & NAND & Logika     & Bitová negácia logického súčinu (NAND) \\
16 & 10000 & NOR  & Logika     & Bitová negácia logického súčtu (NOR) \\
17 & 10001 & NOT  & Logika     & Bitová inverzia operandu (jednotkový doplnok) \\
18 & 10010 & XOR  & Logika     & Bitová neverzia (XOR) \\
19 & 10011 & XNOR & Logika     & Bitová ekvivalencia \\
20 & 10100 & INC  & Aritmetika & Hardvérová inkrementácia operandu o 1 \\
21 & 10101 & DEC  & Aritmetika & Hardvérová dekrementácia operandu o 1 \\
22 & 10110 & PUSH & Zásobník   & Uloženie slova na vrchol zásobníka \\
23 & 10111 & POP  & Zásobník   & Výber slova z vrcholu zásobníka do registra \\
24 & 11000 & BSL  & Posuny     & Barelový posun doľava s konfigurovateľným vstupom \\
25 & 11001 & BSR  & Posuny     & Barelový posun doprava s ukladaním do DBIT \\
26 & 11010 & MSP  & Systém     & Manipulácia so špeciálnymi registrami (FREG, PTREG) \\
27 & 11011 & RSV27 & Rezerva   & Rezervované pre vektorové inštrukcie \\
28 & 11100 & RSV28 & Rezerva   & Rezervované pre rozšírené adresovanie \\
29 & 11101 & RSV29 & Rezerva   & Rezervované pre koprocesor s pohyblivou rádovou čiarkou \\
30 & 11110 & RSV30 & Rezerva   & Rezervované pre kryptografické inštrukcie \\
31 & 11111 & RSV31 & Rezerva   & Rezervované pre budúce rozšírenia architektúry \\
\bottomrule
\end{longtable}
\normalsize

\section{Detailná špecifikácia inštrukčného súboru}
Nasledujúce strany obsahujú úplné architektonické profily všetkých 32 operačných kódov procesora LIM M2.

\InstructionPage{NOP}{00000}{Riadenie}{Hardvérová prázdna operácia a vyprázdnenie aritmetickej reťaze}{%
\small
Signatúra: NOP \\
Operácia:  PC <- PC + 1, ALU\_RESIDUAL <- 0, FREG.CF <- 0, FREG.OF <- 0 \\[2pt]
Príklad:   NOP              ; Reset reťazca prenosu ALU
}{Bez operandov; pole \field{cfg} rezervované na kalibráciu časovania.}{Vynuluje príznaky prenosu (CF), pretečenia (OF) a výpožičky; ostatné príznaky zachováva.}{Slúži na synchronizáciu taktov a hardvérové vyprázdnenie (Flush) vnútorného zvyškového registra ALU, čím prerušuje reťazenie pred novými nezávislými výpočtami.}
\InstructionPage{ADD}{00001}{Aritmetika}{Sčítanie s automatickým reťazením prenosu}{%
\small
Signatúra: ADD arg0(REG), arg1(REG), arg2(REG) : arg0 = arg1 + arg2 \\
Operácia:  sum <- arg1 + arg2 + ALU\_RESIDUAL \\
           arg0 <- sum[15:0], ALU\_RESIDUAL <- (sum > 0xFFFF ? 1 : 0) \\[2pt]
Príklad:   ADD R0, R1, R2   ; R0 = R1 + R2 (registre R0--R7)
}{Registre R0--R7, akumulátory ACCA/ACCB alebo priame konštanty.}{Nastavuje príznak prenosu CF a pretečenia OF pri prekročení $2^{16}-1$; aktualizuje ZF podľa dst.}{Podporuje priebežnú viacbajtovú aritmetiku (32, 48, 64+ bitov). Ak súčet prekročí 16 bitov ($>2^{16}-1$), prenos (+1) sa zachová vo vnútornom registri ALU a automaticky sa pripočíta pri nasledujúcej operácii ADD bez nutnosti inštrukcie ADC. Spodných 16 bitov sa deterministicky zapíše do dst. Reťazenie sa preruší inštrukciou NOP alebo novou matematickou operáciou.}
\InstructionPage{SUB}{00010}{Aritmetika}{Odčítanie s automatickým reťazením výpožičky}{%
\small
Signatúra: SUB arg0(REG), arg1(REG), arg2(REG) : arg0 = arg1 - arg2 \\
Operácia:  diff <- arg1 - arg2 - ALU\_RESIDUAL \\
           arg0 <- diff[15:0], ALU\_RESIDUAL <- (diff < 0 ? 1 : 0) \\[2pt]
Príklad:   SUB R0, R1, R2   ; R0 = R1 - R2 (registre R0--R7)
}{Registre R0--R7, konštanty alebo akumulátory ACCA/ACCB.}{Pri zápornom rozdiele ($< 0$) nastavuje príznak výpožičky (Borrow / CF) a OF; aktualizuje ZF.}{Umožňuje kontinuálne viacbajtové odčítanie. Pri zápornom rozdiele ($< 0$) sa výpožička ($-1$) uloží do vnútorného registra ALU a pri ďalšej inštrukcii SUB sa automaticky odpočíta 1 od vyššieho slova. Výsledok spodného slova sa deterministicky uloží do dst. Na začatie nového nezávislého odčítania sa reťazenie preruší pomocou NOP.}
\InstructionPage{MUL}{00011}{Aritmetika}{Celočíselné násobenie so zachovaním hornej časti v registri ALU}{%
\small
Signatúra: MUL arg0(REG), arg1(REG), arg2(REG) : arg0 = arg1 * arg2 \\
Operácia:  prod <- arg1 * arg2 \\
           arg0 <- prod[15:0], ALU\_RESIDUAL <- prod[31:16] \\[2pt]
Príklad:   MUL R0, R1, R2   ; R0 = R1 * R2 (horná časť v ALU\_RESIDUAL)
}{Registre R0--R7, akumulátory ACCA/ACCB.}{Aktualizuje ZF; nastavuje OF/CF, ak je horných 16 bitov súčinu (prod[31:16]) nenulových.}{Pri 32-bitovom súčine smeruje spodných 16 bitov do cieľového registra dst, kým horných 16 bitov sa uloží do vnútorného zvyškového registra ALU (ALU Residual). Pri nasledujúcom sčítaní sa tento zvyšok automaticky pripočíta (podpora MAC a viacbajtových čísiel). Hornú polovicu možno získať aj sčítaním registra so sebou samým po vynulovaní. Inštrukcia NOP register vyčistí.}
\InstructionPage{DIV}{00100}{Aritmetika}{Celočíselné delenie so zachovaním zvyšku vo vnútornom registri ALU}{%
\small
Signatúra: DIV arg0(REG), arg1(REG), arg2(REG) : arg0 = arg1 / arg2 \\
Operácia:  arg0 <- arg1 / arg2 \\
           ALU\_RESIDUAL <- arg1 \% arg2 \\[2pt]
Príklad:   DIV R0, R1, R2   ; R0 = R1 / R2 (zvyšok modulo v ALU\_RESIDUAL)
}{Registre R0--R7, akumulátory ACCA/ACCB.}{Pri delení nulou sa nastavuje príznak chyby EM (Error Math) bez zablokovania procesora; aktualizuje ZF.}{Podiel sa zapíše do cieľového registra dst, pričom zvyšok po delení (modulo) sa zachová vo vnútornom registri ALU pre ďalšie operácie alebo vymazanie inštrukciou NOP.}
\InstructionPage{LMR}{00101}{Pamäť}{Načítanie dát z pamäte a externých zariadení do registrov}{%
\small
Signatúra: LMR arg0(REG), <SRC\_MODE> : arg0 = MEM[...] \\
Režimy:    SRC=0: LMR R0, [PC+]          ; Imm Data (1--2 bajty) \\
           SRC=1: LMR R0, [[PC+]]        ; Imm Ukazovateľ \\
           SRC=2: LMR R0, [PTREG]        ; Systémový ukazovateľ (AIP) \\
           SRC=3: LMR R0(R\_hi), R1(R\_lo) ; 24-bitová adresa z páru \\
Operácia:  arg0 <- MEM[EA] (horné bity adresy sa prepíšu dátami) \\[2pt]
Príklad:   LMR R0, R1       ; EA=(R0<<16)|R1, načítanie MEM[EA] do R0
}{%
\textbf{Oktet 0:} bit $[8]$ (\field{SZ}) určuje šírku dát ($0 = 1$ bajt, $1 = 2$ bajty / slovo); bity $[10:9]$ (\field{SRC}) volia zdroj dát ($0$ --- priame dáta za kódom inštrukcie, $1$ --- priamy ukazovateľ za inštrukciou, $2$ --- nepriamo cez 24-bitový ukazovateľ \reg{PTREG}, $3$ --- nepriamo cez dvojicu používateľských registrov).\\
\textbf{Oktet 1:} bity $[7:5]$ definujú register \reg{Arg1} (\reg{R\_hi} / cieľ), bity $[4:2]$ definujú register \reg{Arg2} (\reg{R\_lo}). V režime $\field{SRC} = 3$ používateľ zvolí dva 16-bitové registre tvoriace 32-bitové číslo, z ktorého sa na adresovanie použije 24 bitov: \reg{Arg1} nesie horné bity adresy a \reg{Arg2} dolných 16 bitov. Register \reg{Arg1} je zároveň cieľovým registrom: po výbere z pamäte sa horné bity adresy v \reg{Arg1} prepíšu načítanými dátami.%
}{Všetky príznaky v stavovom registri \reg{FREG} zostávajú zachované bez zmeny.}{%
Inštrukcia \texttt{LMR} predstavuje kľúčovú vstupnú bránu pre prenos externých dát a pamäte do registrového poľa procesora LIM M2 ($R0$--$R7$). Architektúra integruje štyri ortogonálne režimy výberu s natívnym 24-bitovým adresovaním:\\
1. V režime $\field{SRC} = 0$ procesor číta literál priamo z toku kódu ($1$ alebo $2$ bajty) s automatickým posunom čítača inštrukcií \reg{PC}.\\
2. V režime $\field{SRC} = 1$ sa za inštrukciou načíta absolútny ukazovateľ, podľa ktorého sa vyberú dáta z pamäte.\\
3. V režime $\field{SRC} = 2$ prebieha prístup na 24-bitovú fyzickú adresu z registra \reg{PTREG}. Pri aktívnom príznaku autoincrementu \texttt{FREG.AIP} (bit 12) sa register \reg{PTREG} automaticky zvýši o veľkosť výberu ($+1$ alebo $+2$), čo umožňuje prúdové načítanie polí za jediný takt.\\
4. V režime $\field{SRC} = 3$ používateľ zvolí dva 16-bitové registre (\reg{Arg1} a \reg{Arg2}), ktoré spolu vytvárajú 32-bitové číslo, z ktorého sa v praxi adresovania využíva dolných 24 bitov. Používateľ určí, že horné bity adresy sú v \reg{Arg1} a dolných 16 bitov v \reg{Arg2}. V tomto režime sa register \reg{Arg1} s hornými bitmi adresy stáva zároveň cieľovým registrom, do ktorého sa zapíšu načítané dáta (teda horné bity adresy v \reg{Arg1} sa prepíšu dátami, kým \reg{Arg2} uchováva dolných 16 bitov adresy).%
}
\InstructionPage{LRM}{00110}{Pamäť}{Zápis dát z registrov do pamäte a externých zariadení}{%
\small
Signatúra: LRM arg0(REG), <DST\_MODE> : MEM[...] = arg0 \\
Režimy:    DST=0: LRM R0, [PC+]          ; Vložená pamäť za kódom \\
           DST=1: LRM R0, [[PC+]]        ; Imm Ukazovateľ \\
           DST=2: LRM R0, [PTREG]        ; Systémový ukazovateľ (AIP) \\
           DST=3: LRM R0(R\_hi), R1(R\_lo) ; 24-bitová adresa z páru \\
Operácia:  MEM[EA] <- arg0[15:0] (alebo dolný bajt pri SZ=0) \\[2pt]
Príklad:   LRM R0, R1       ; EA=(R0<<16)|R1, zápis R0 do MEM[EA]
}{%
\textbf{Oktet 0:} bit $[8]$ (\field{SZ}) určuje šírku dát ($0 = 1$ bajt, $1 = 2$ bajty / slovo); bity $[10:9]$ (\field{DST}) volia cieľové umiestnenie dát ($0$ --- zápis priamo za kód inštrukcie, $1$ --- zápis podľa ukazovateľa za inštrukciou, $2$ --- nepriamo cez 24-bitový ukazovateľ \reg{PTREG}, $3$ --- nepriamo cez dvojicu registrov používateľa).\\
\textbf{Oktet 1:} bity $[7:5]$ definujú zdrojový register \reg{Arg1} (\reg{Rs} / \reg{R\_hi}), bity $[4:2]$ definujú register \reg{Arg2} (\reg{R\_lo}). V režime $\field{DST} = 3$ dva registre tvoria 32-bitové číslo, z ktorého sa použije 24 bitov adresy (\reg{Arg1} --- horné bity adresy a zdroj dát, \reg{Arg2} --- dolných 16 bitov). Prepísanie existujúcich dát na cieľovej adrese je plnou zodpovednosťou programátora.%
}{Všetky príznaky v stavovom registri \reg{FREG} zostávajú zachované bez zmeny.}{%
Inštrukcia \texttt{LRM} predstavuje symetrický mechanizmus prenosu registrových dát procesora LIM M2 ($R0$--$R7$) do pamäte a periférnych obvodov. Architektúra integruje štyri ortogonálne režimy adresovania:\\
1. V režime $\field{DST} = 0$ procesor zapisuje $1$ alebo $2$ bajty priamo do pamäte za inštrukciou (\reg{PC}) s posunom čítača inštrukcií. Režim slúži pre dynamické buffery a samomodifikujúci sa kód; prepísanie dát je na zodpovednosti programátora.\\
2. V režime $\field{DST} = 1$ sa za inštrukciou načíta 16-bitový absolútny ukazovateľ, na ktorý sa dáta uložia.\\
3. V režime $\field{DST} = 2$ zápis smeruje na 24-bitovú systémovú adresu registra \reg{PTREG}. Pri aktívnom príznaku \texttt{FREG.AIP} (bit 12) sa register \reg{PTREG} automaticky zvýši o $+1$ alebo $+2$, čo umožňuje prúdový DMA výstup do pamäte za jediný takt.\\
4. V režime $\field{DST} = 3$ sa využíva dvojica 16-bitových registrov: \reg{Arg1} a \reg{Arg2} tvoria 32-bitové číslo, ktorého dolných 24 bitov určuje fyzickú adresu ($\text{Addr}[23:0]$). Dáta z \reg{Arg1} sa zapíšu na vypočítanú adresu, čo sprístupňuje celých 16 MB priestoru.%
}
\InstructionPage{LRR}{00111}{Presun}{Medziregistrový presun (kopírovanie z bodu B do bodu A)}{%
\small
Signatúra: LRR arg0(REG), arg1(REG) : arg0 = arg1 \\
Operácia:  arg0 <- arg1 (priame kopírovanie z bodu B do bodu A) \\[2pt]
Príklad:   LRR R0, R1       ; R0 = R1 (hodnota R1 sa skopíruje do R0)
}{%
\textbf{Oktet 0:} kód operácie $00111_2$ (\texttt{0x07}), bity $[10:8]$ sú rezervované ($000_2$).\\
\textbf{Oktet 1:} bity $[7:5]$ definujú cieľový register-príjemcu \reg{Arg1} (\reg{Rd}, bod A, kam sa dáta ukladajú); bity $[4:2]$ definujú zdrojový register \reg{Arg2} (\reg{Rs}, bod B, odkiaľ sa dáta čítajú); bity $[1:0]$ obsahujú $00_2$. Pri vykonaní sa hodnota \reg{Rs} prenesie do \reg{Rd} za 1 mikrocyklus.%
}{Všetky príznaky zostávajú zachované.}{%
Inštrukcia \texttt{LRR} predstavuje základnú operáciu presunu dát medzi registrami všeobecného určenia procesora ($R0$--$R7$). Vykonáva prenos zo zdroja (bod B) do cieľa (bod A). Presun sa uskutočňuje cez priamy hardvérový premosťovací kanál \textbf{MUX 1} $\to$ \textbf{MUX 2} bez zaťažovania ALU a bez vplyvu na stav príznakov registra \reg{FREG}.%
}
\InstructionPage{CMP}{01000}{Porovnanie}{Hardvérové porovnanie operandov}{%
\small
Signatúra: CMP arg0(REG), arg1(REG) : compare(arg0, arg1) \\
Operácia:  if arg0 < arg1 then FREG.CMP <- 01b (menšie) \\
           if arg0 == arg1 then FREG.CMP <- 10b (rovné) \\
           if arg0 > arg1 then FREG.CMP <- 11b (väčšie) \\
           (stav 00b --- počiatočný reset / neporovnané / chyba) \\[2pt]
Príklad:   CMP R0, R1       ; 01: R0<R1, 10: R0==R1, 11: R0>R1 (00: reset/chyba)
}{Režim porovnania registrov bez zmeny dát.}{Aktualizuje 2-bitové pole FREG.CMP (bity [5:4]) a príznaky ZF, CF, OF bez zmeny registrov.}{Umožňuje rýchle podmienené vetvenie.}
\InstructionPage{JMP}{01001}{Skok}{Nepodmienený skok}{%
\small
Signatúra: JMP target(IMM|REG) : PC = target \\
Operácia:  PC <- target \\[2pt]
Príklad:   JMP R0           ; Bezpodmienečný skok na adresu v R0
}{Podporuje 16-bitový lokálny alebo plný 24-bitový skok cez \reg{PTREG}.}{Príznaky zostávajú nezmenené.}{Atomické načítanie Program Counter v minimálnom počte cyklov.}
\InstructionPage{JCR}{01010}{Skok}{Podmienený skok pri prenose}{%
\small
Signatúra: JCR target(IMM|REG) : if FREG.CF == 1 then PC = target \\
Operácia:  if FREG.CF == 1 then PC <- target \\[2pt]
Príklad:   JCR R2           ; Skok na adresu R2, ak je nastavený Carry
}{Cieľová adresa určená posunom alebo absolútne.}{Príznaky sú iba čítané.}{Základný stavebný prvok viacnásobnej presnej aritmetiky.}
\InstructionPage{JNZ}{01011}{Skok}{Podmienený skok pri nenulovom stave}{%
\small
Signatúra: JNZ target(IMM|REG) : if FREG.ZF == 0 then PC = target \\
Operácia:  if FREG.ZF == 0 then PC <- target \\[2pt]
Príklad:   JNZ R3           ; Skok na adresu R3, ak ZF == 0 (nie je nula)
}{Vysoko efektívny skok pre programové slučky a počítadlá.}{Príznaky zostávajú nezmenené.}{Rýchle vyhodnotenie podmienky v mikroprograme.}
\InstructionPage{JZ}{01100}{Skok}{Podmienený skok pri nulovom stave}{%
\small
Signatúra: JZ target(IMM|REG) : if FREG.ZF == 1 then PC = target \\
Operácia:  if FREG.ZF == 1 then PC <- target \\[2pt]
Príklad:   JZ R4            ; Skok na adresu R4, ak ZF == 1 (je nula)
}{Vyhodnocuje rovnosť alebo ukončovacie podmienky.}{Príznaky zostávajú nezmenené.}{Optimalizované pre reťazcové a vyhľadávacie operácie.}
\InstructionPage{AND}{01101}{Logika}{Bitový logický súčin (AND)}{%
\small
Signatúra: AND arg0(REG), arg1(REG) : arg0 = arg0 \& arg1 \\
Operácia:  arg0 <- arg0 AND arg1 \\[2pt]
Príklad:   AND R0, R1   ; R0 = R0 \& R1
}{Registrové operandy (2 operandy).}{Nastavuje príznak nuly ZF; pretečenie vynulované.}{Kľúčová inštrukcia pre bitové maskovanie.}
\InstructionPage{OR}{01110}{Logika}{Bitový logický súčet (OR)}{%
\small
Signatúra: OR arg0(REG), arg1(REG) : arg0 = arg0 | arg1 \\
Operácia:  arg0 <- arg0 OR arg1 \\[2pt]
Príklad:   OR R0, R1    ; R0 = R0 | R1
}{Registrové operandy (2 operandy).}{Nastavuje ZF; prenosy sú deterministicky vynulované.}{Rýchle zlučovanie riadiacich príznakov.}
\InstructionPage{NAND}{01111}{Logika}{Bitový logický zápor súčinu (NAND)}{%
\small
Signatúra: NAND arg0(REG), arg1(REG) : arg0 = \textasciitilde(arg0 \& arg1) \\
Operácia:  arg0 <- NOT (arg0 AND arg1) \\[2pt]
Príklad:   NAND R0, R1  ; R0 = \textasciitilde(R0 \& R1)
}{Univerzálna logická báza (2 operandy).}{Nastavuje príznak ZF.}{Umožňuje kompaktnú realizáciu zákazníckych logických funkcií.}
\InstructionPage{NOR}{10000}{Logika}{Bitový logický zápor súčtu (NOR)}{%
\small
Signatúra: NOR arg0(REG), arg1(REG) : arg0 = \textasciitilde(arg0 | arg1) \\
Operácia:  arg0 <- NOT (arg0 OR arg1) \\[2pt]
Príklad:   NOR R0, R1   ; R0 = \textasciitilde(R0 | R1)
}{Invertovaná logická funkcia (2 operandy).}{Nastavuje príznak ZF.}{Rozširuje množinu rýchlych bitových transformácií.}
\InstructionPage{NOT}{10001}{Logika}{Bitová inverzia (jednotkový doplnok)}{%
\small
Signatúra: NOT arg0(REG) : arg0 = \textasciitilde arg0 \\
Operácia:  arg0 <- NOT arg0 \\[2pt]
Príklad:   NOT R0       ; R0 = \textasciitilde R0 (1-bajtová inštrukcia)
}{Jednobajtová inštrukcia ([7:3] opkód 10001, [2:0] Rd). Druhý bajt z pamäte sa nenačítava.}{Nastavuje príznak ZF.}{Vykoná sa v minimálnom počte mikrotaktov bez načítania druhého bajtu.}
\InstructionPage{XOR}{10010}{Logika}{Bitová neverzia (XOR)}{%
\small
Signatúra: XOR arg0(REG), arg1(REG) : arg0 = arg0 \textasciicircum\ arg1 \\
Operácia:  arg0 <- arg0 XOR arg1 \\[2pt]
Príklad:   XOR R0, R1   ; R0 = R0 \textasciicircum\ R1 (pri R0==R1 vynuluje R0)
}{Registrové operandy (2 operandy).}{Nastavuje príznak ZF.}{Základná inštrukcia pre paritu, CRC a kryptografiu.}
\InstructionPage{XNOR}{10011}{Logika}{Bitová ekvivalencia}{%
\small
Signatúra: XNOR arg0(REG), arg1(REG) : arg0 = \textasciitilde(arg0 \textasciicircum\ arg1) \\
Operácia:  arg0 <- NOT (arg0 XOR arg1) \\[2pt]
Príklad:   XNOR R0, R1  ; R0 = \textasciitilde(R0 \textasciicircum\ R1)
}{Porovnanie bitov na zhodu (2 operandy).}{Nastavuje príznak ZF.}{Rýchla paralelná kontrola ekvivalencie masiek.}
\InstructionPage{INC}{10100}{Aritmetika}{Inkrementácia o 1}{%
\small
Signatúra: INC arg0(REG) : arg0 = arg0 + 1 \\
Operácia:  arg0 <- arg0 + 1 \\[2pt]
Príklad:   INC R0           ; R0 = R0 + 1 (inkrement registra)
}{Rýchly krok počítadla cyklu.}{Nastavuje ZF, OF; prenos CF riadený cez \field{cfg}.}{Optimalizuje indexovanie v poliach.}
\InstructionPage{DEC}{10101}{Aritmetika}{Dekrementácia o 1}{%
\small
Signatúra: DEC arg0(REG) : arg0 = arg0 - 1 \\
Operácia:  arg0 <- arg0 - 1 \\[2pt]
Príklad:   DEC R0           ; R0 = R0 - 1 (dekrement registra)
}{Krok zostupného počítadla.}{Nastavuje ZF, OF.}{Kompaktná realizácia cyklov.}
\InstructionPage{PUSH}{10110}{Zásobník}{Uloženie slova do zásobníka}{%
\small
Signatúra: PUSH arg0(REG) : SREG -= 2, MEM[SREG] = arg0 \\
Operácia:  SREG <- SREG - 2, MEM[SREG] <- arg0 \\[2pt]
Príklad:   PUSH R0          ; Uložiť slovo R0 na vrchol zásobníka
}{Zdroj: ľubovoľný register R0--R7 alebo systémový register.}{Aritmetické príznaky sa nemenia.}{Hardvérová podpora volaní funkcií a ukladania kontextu.}
\InstructionPage{POP}{10111}{Zásobník}{Výber slova zo zásobníka}{%
\small
Signatúra: POP arg0(REG) : arg0 = MEM[SREG], SREG += 2 \\
Operácia:  arg0 <- MEM[SREG], SREG <- SREG + 2 \\[2pt]
Príklad:   POP R0           ; Vybrať slovo z vrcholu zásobníka do R0
}{Cieľ: registre R0--R7 alebo stavové registre.}{Aritmetické príznaky sa nemenia.}{Symetrické obnovenie kontextu po prerušení.}
\InstructionPage{BSL}{11000}{Posuny}{Barelový posun doľava}{%
\small
Signatúra: BSL arg0(REG), arg1(REG), arg2(IMM|REG) : arg0 = arg1 << arg2 \\
Operácia:  arg0 <- arg1 << arg2 \\
           FREG.DBIT <- vytlačený najvyšší bit (bit 15) \\
           Pri SMOD=1 sa do LSB vloží predošlý DBIT; pri SMOD=0: 0 \\[2pt]
Príklad:   BSL R0, R1, 1    ; Posun R1 doľava o 1 bit, výsledok do R0
}{Rozsah posunu: konštanta alebo hodnota v registri.}{Najvyšší vysunutý bit zachytený v FREG.DBIT (bit 13). Pri SMOD=1 sa do najnižšieho bitu vloží predchádzajúci DBIT; pri SMOD=0 sa dopĺňajú nuly.}{Efektívna cyklická a prúdová rotácia viacbajtových dát.}
\InstructionPage{BSR}{11001}{Posuny}{Barelový posun doprava}{%
\small
Signatúra: BSR arg0(REG), arg1(REG), arg2(IMM|REG) : arg0 = arg1 >> arg2 \\
Operácia:  arg0 <- arg1 >> arg2 \\
           FREG.DBIT <- vytlačený najnižší bit (bit 0) \\
           Pri SMOD=1 sa do MSB vloží predošlý DBIT; pri SMOD=0: 0 \\[2pt]
Príklad:   BSR R0, R1, 1    ; Posun R1 doprava o 1 bit, výsledok do R0
}{Logický alebo aritmetický posun voliteľný cez \field{cfg}.}{Najnižší vysunutý bit uložený v DBIT (bit 13). Pri SMOD=1 sa do najvyššieho bitu vloží predchádzajúci DBIT; pri SMOD=0 sa dopĺňajú nuly.}{Prúdové škálovanie a delenie mocninami dvojky.}
\InstructionPage{MSP}{11010}{Systém}{Manipulácia so špeciálnymi registrami (FREG, PTREG)}{%
\small
Signatúra: MSP sreg, [op], args : práca so špeciálnymi registrami \\
Operácia:  
  \textbf{Priamy zápis (bit 10 = 0):} \\
  \quad sreg=00 (FREG):  FREG <- Rs[15:0] (oktet 1: [7:5]=Rs) \\
  \quad sreg=01 (PTREG): PTREG <- (Rs\_hi $\ll$ 16) | Rs\_lo (oktet 1: [7:5]=Rs\_hi, [4:2]=Rs\_lo) \\
  \textbf{Modifikácia (bit 10 = 1):} \\
  \quad sreg=00 (FREG):  FREG <- FREG $\odot$ Rs (oktet 1: [7:6]=op \{AND,OR,XOR,XNOR\}, [5:3]=Rs) \\
  \quad sreg=01 (PTREG): PTREG <- PTREG $\pm$ Rs (oktet 1: [7]=op \{ADD,SUB\}, [6:4]=Rs) \\[2pt]
Príklad:   MSP FREG, R0         ; Priamy zápis do FREG hodnoty R0 \\
           MSP PTREG, R0, R1    ; Nastavenie 24-bitového PTREG = \{R0, R1\} \\
           MSP FREG, OR, R2     ; Nastavenie masiek: FREG = FREG | R2 \\
           MSP PTREG, ADD, R3   ; Posun ukazovateľa: PTREG = PTREG + R3
}{Špeciálne registre FREG (00) a 24-bitový PTREG (01). Podporuje priamu inicializáciu, bitové masky FREG a lokálnu aritmetiku PTREG.}{Príznaky FREG sa modifikujú priamo alebo maskou. Aritmetika PTREG generuje príznak prenosu CF cez ALU.}{Univerzálny privilegovaný riadiaci primitív systémovej úrovne.}
\InstructionPage{RSV27}{11011}{Rezerva}{Hardvérové rozšírenie}{%
\small
Signatúra: RSV27 : rezervované \\
Operácia:  Reserved for LIM architecture hardware extensions
}{Vyhradené pre budúce revízie architektúry LIM.}{Nedefinované.}{Rezervované.}
\InstructionPage{RSV28}{11100}{Rezerva}{Hardvérové rozšírenie}{%
\small
Signatúra: RSV28 : rezervované \\
Operácia:  Reserved for LIM architecture hardware extensions
}{Vyhradené pre budúce revízie architektúry LIM.}{Nedefinované.}{Rezervované.}
\InstructionPage{RSV29}{11101}{Rezerva}{Hardvérové rozšírenie}{%
\small
Signatúra: RSV29 : rezervované \\
Operácia:  Reserved for LIM architecture hardware extensions
}{Vyhradené pre budúce revízie architektúry LIM.}{Nedefinované.}{Rezervované.}
\InstructionPage{RSV30}{11110}{Rezerva}{Hardvérové rozšírenie}{%
\small
Signatúra: RSV30 : rezervované \\
Operácia:  Reserved for LIM architecture hardware extensions
}{Vyhradené pre budúce revízie architektúry LIM.}{Nedefinované.}{Rezervované.}
\InstructionPage{RSV31}{11111}{Rezerva}{Hardvérové rozšírenie}{%
\small
Signatúra: RSV31 : rezervované \\
Operácia:  Reserved for LIM architecture hardware extensions
}{Vyhradené pre budúce revízie architektúry LIM.}{Nedefinované.}{Rezervované.}

\clearpage
\section{Fyzické rozhranie a rozloženie vývodov (34-vývodové puzdro)}
Procesor je integrovaný v kompaktnom 34-vývodovom puzdre s dvojradovým usporiadaním kontaktov:

\begin{center}
\begin{tikzpicture}[font=\sffamily\scriptsize,
  pinbox/.style={draw,thick,fill=white,minimum width=6mm,minimum height=4mm,inner sep=0pt,font=\bfseries\tiny}]
\def\chipW{3.2}
\def\chipH{9.9}
\draw[thick,fill=limlight] (0,0) rectangle (\chipW,-\chipH);
\draw[thick] ($(\chipW/2,0)+(-0.35,0)$) arc (180:360:0.35);
\node[font=\bfseries\Large,text=limblue,rotate=-90] at (\chipW/2,-\chipH/2) {LIM M2};
\foreach \i/\sigL/\sigR in {
  1/A0/D0, 2/A1/D1, 3/A2/D2, 4/A3/D3, 5/A4/D4, 6/A5/D5, 7/A6/D6, 8/A7/D7,
  9/A8/IA, 10/A9/VCC, 11/A10/GND, 12/A11/RD, 13/A12/WD, 14/A13/RF, 15/A14/AL, 16/A15/RDI, 17/CLK1/CLK2
} {
  \pgfmathsetmacro{\y}{-\i * 0.55}
  \pgfmathtruncatemacro{\pinR}{35 - \i}
  \node[pinbox] at (0,\y) {\i};
  \node[anchor=east] at (-0.38,\y) {\sigL};
  \node[pinbox] at (\chipW,\y) {\pinR};
  \node[anchor=west] at (\chipW+0.38,\y) {\sigR};
}
\end{tikzpicture}
\end{center}

\subsection{Funkčný popis signálov}
\begin{tabularx}{\textwidth}{>{\ttfamily}p{28mm}p{24mm}X}
\toprule
Vývod & Typ & Popis systémového signálu \\
\midrule
A0--A15    & Výstup / Tri-State & 16-bitová adresová zbernica (riadená externým impulzom AL). \\
D0--D7     & Obojsmerný         & 8-bitová zbernica pre prenos dát s pamäťou a perifériami. \\
CLK1, CLK2 & Vstup              & Dvojfázové hodinové signály s ochranou proti prekrytiu fáz. \\
VCC, GND   & Napájanie          & Hlavné napájacie vedenie logiky procesora (+5 V nominálne). \\
RD, WD     & Výstup             & Riadiace impulzy čítania (Read) a zápisu (Write) na zbernici. \\
AL         & Výstup             & Address Latch: riadenie záchytného registra adresy. \\
RF         & Vstup              & Ready Flag: hardvérové potvrdenie pripravenosti podriadeného zariadenia. \\
RDI        & Vstup              & Raw Direct Interrupt: prioritná linka hardvérového prerušenia. \\
IA         & Výstup             & Interruption Accepted / Interrupt Acknowledge: potvrdenie prijatia prerušenia; procesor očakáva načítanie čísla vektora zo zbernice dát. \\
\bottomrule
\end{tabularx}

\clearpage
\section{Podsystém prerušení a vektorové spracovanie}

Podsystém prerušení procesora LIM~M2 je navrhnutý tak, aby poskytoval minimálnu latenciu odozvy na udalosti v reálnom čase, prísny determinizmus a úplnú ochranu pred zbernicovými kolíziami. Architektúra zjednocuje externé hardvérové požiadavky a interné softvérové pasce do jednotnej 16-vektorovej hierarchie s hardvérovým potvrdzovaním na fyzickej úrovni.

\subsection{Zjednotený priestor 16 vektorov prerušení}
Procesor podporuje presne 16 systémových vektorov (indexy 0--15), ktoré sú ortogonálne zdieľané medzi externými perifériami a internými udalosťami procesorového jadra:

\begin{table}[h]
\centering
\caption{Zjednotená tabuľka 16 vektorov prerušení LIM M2}
\label{tab:vectors}
\begin{tabularx}{\textwidth}{>{\bfseries}c >{\ttfamily}c p{32mm}X}
\toprule
Vektor & Kód & Zdroj / Typ & Funkčné určenie \\
\midrule
VEC 0  & 0x00 & Hardvérový / Reset  & Power-on Reset: počiatočná inicializácia jadra a registrov. \\
VEC 1  & 0x01 & Hardvérový (NMI)    & Nemaskovateľné prerušenie: výpadok napájania, kritická porucha. \\
VEC 2  & 0x02 & Softvérová pasca    & Aritmetická výnimka (delenie nulou, príznak \reg{FREG.EM}). \\
VEC 3  & 0x03 & Softvérová pasca    & Systémové volanie jadra (Syscall / API dispečer OS). \\
VEC 4  & 0x04 & Hardvérový / Perif. & Systémový intervalový časovač / Kvantovač plánovača úloh. \\
VEC 5  & 0x05 & Hardvérový / Perif. & Sériové komunikačné rozhranie (UART RX/TX). \\
VEC 6  & 0x06 & Hardvérový / Perif. & Radič priameho prístupu do pamäte (DMA Transfer Complete). \\
VEC 7  & 0x07 & Hardvérový / Perif. & Externý radič prerušení / Zbernicový most. \\
VEC 8--15 & 0x08--0x0F & Zdieľané (HW / SW) & Používateľské hardvérové linky a softvérové vektory. \\
\bottomrule
\end{tabularx}
\end{table}

\subsection{Hardvérový protokol potvrdzovania (RDI $\to$ IA $\to$ D0--D7)}
Pre spoľahlivú komunikáciu s externými radičmi je implementovaný 4-fázový cyklus hardvérového potvrdzovania (handshake), ktorý úplne eliminuje réžiu cyklického dopytovania (polling):

\begin{enumerate}
  \item \textbf{Fáza požiadavky (RDI):} Periférne zariadenie nastaví aktívnu vysokú úroveň na vstupe \texttt{RDI} (\textit{Raw Direct Interrupt}). Linka je testovaná jadrom na atomických hraniciach mikroinštrukcií.
  \item \textbf{Fáza kontroly masky (Arbitrácia):} Ak je globálne maskovanie prerušení povolené (bit \reg{FREG.MIE} = 1), procesor požiadavku zachytí a dokončí aktuálnu mikrooperáciu dátovej cesty.
  \item \textbf{Fáza potvrdenia (Strob IA):} Jadro vygeneruje aktívny strob na výstupe \texttt{IA} (\textit{Interruption Accepted / Interrupt Acknowledge}, pin~26). Tento strob informuje všetky periférie na zbernici o pripravenosti jadra prijať číslo vektora.
  \item \textbf{Fáza prenosu vektora (Zbernica D0--D7):} Súčasne s aktiváciou \texttt{IA} sa dátová zbernica \texttt{D0--D7} prepne do režimu vstupu a aktivuje sa signál čítania \texttt{RD}. Žiadajúce zariadenie vystaví na linky \texttt{D0--D7} 8-bitový identifikátor vektora (spodné 4 bity adresujú vektory 0--15).
  \item \textbf{Fáza záchytu a prechod:} Procesor zachytí dáta zo zbernice \texttt{D0--D7}, deaktivuje signál \texttt{IA} a prejde k uloženiu kontextu.
\end{enumerate}

\begin{center}
\begin{tikzpicture}[x=1.3cm,y=0.65cm,>=Stealth,font=\sffamily\scriptsize]
  \foreach \t in {0,...,6} {
    \draw[dotted,gray!50] (\t, 0.5) -- (\t, 5.5);
    \node[above,gray!70,font=\tiny] at (\t+0.5, 5.5) {T\t};
  }
  
  \node[left,font=\bfseries] at (0, 5) {CLK};
  \draw[thick,blue!80!black] (0, 4.7) 
    \foreach \t in {0,...,5} { -- ++(0.5, 0.6) -- ++(0.5, 0) -- ++(0, -0.6) -- ++(0, 0) };

  \node[left,font=\bfseries] at (0, 4) {RDI};
  \draw[thick,red!80!black] (0, 3.8) -- (1.0, 3.8) -- (1.2, 4.4) -- (4.8, 4.4) -- (5.0, 3.8) -- (6.0, 3.8);
  \node[above,red!80!black,font=\tiny] at (2.0, 4.4) {Požiadavka zariadenia};

  \node[left,font=\bfseries] at (0, 3) {IA};
  \draw[thick,teal!80!black] (0, 2.8) -- (2.0, 2.8) -- (2.2, 3.4) -- (4.2, 3.4) -- (4.4, 2.8) -- (6.0, 2.8);
  \node[above,teal!80!black,font=\tiny] at (3.2, 3.4) {Strob IA (Ack)};

  \node[left,font=\bfseries] at (0, 2) {D0--D7};
  \draw[thick,gray!60] (0, 1.8) -- (2.2, 1.8);
  \draw[thick,black] (2.2, 1.8) -- (2.4, 2.2) -- (4.0, 2.2) -- (4.2, 1.8) -- (6.0, 1.8);
  \draw[thick,black] (2.2, 1.8) -- (2.4, 1.4) -- (4.0, 1.4) -- (4.2, 1.8);
  \node[font=\bfseries\tiny] at (3.2, 1.8) {Vektor [3:0]};

  \node[left,font=\bfseries] at (0, 1) {JADRO};
  \draw[thick,cyan!80!black] (0, 0.8) -- (2.0, 0.8) -- (2.2, 1.2) -- (4.2, 1.2) -- (4.4, 0.8) -- (6.0, 0.8);
  \node[above,cyan!80!black,font=\tiny] at (1.0, 0.8) {Mikrokód};
  \node[above,cyan!80!black,font=\tiny] at (3.2, 1.2) {Príjem vektora};
  \node[above,cyan!80!black,font=\tiny] at (5.2, 0.8) {Zásobník (PC, FREG)};
\end{tikzpicture}
\end{center}

\subsection{Uloženie kontextu v 24-bitovom zásobníku a vstup do obsluhy}
Prijatie prerušenia je povolené výlučne za podmienky, že procesor už nevykonáva iné prerušenie (\reg{FREG.INTR} = 0) a prerušenia sú globálne povolené (\reg{FREG.MIE} = 1). Ak \reg{FREG.INTR} = 1, akékoľvek ďalšie hardvérové prerušenia sú hardvérovo zablokované.

Po potvrdení prerušenia jadro vykoná nasledujúci postup ochrany stavu:
\begin{enumerate}
  \item \textbf{Uloženie čítača inštrukcií (PC):} Aktuálna hodnota \texttt{PC} sa uloží na vrchol systémového zásobníka posunom 24-bitového ukazovateľa \reg{SREG}:
  $$\reg{SREG} \leftarrow \reg{SREG} - 2, \quad \text{Memory}[\reg{SREG}] \leftarrow \text{PC}_{15:0}$$
  \item \textbf{Uloženie príznakov a stavu (FREG):} Aktuálny stav registra \reg{FREG} (s pôvodnou hodnotou \reg{FREG.INTR}=0) sa zapíše do zásobníka:
  $$\reg{SREG} \leftarrow \reg{SREG} - 2, \quad \text{Memory}[\reg{SREG}] \leftarrow \reg{FREG}$$
  \item \textbf{Hardvérové blokovanie vnorených prerušení:} Jadro nastaví príznak vykonávania prerušenia \reg{FREG.INTR} $\leftarrow$ \texttt{1} (bit 15) a vynuluje masku \reg{FREG.MIE} $\leftarrow$ \texttt{0}, čím sa zamedzí kaskádovým reentrantným volaniam.
  \item \textbf{Skok na vektor:} Adresa obslužnej rutiny sa vypočíta z prijatého čísla vektora a zapíše do \texttt{PC}.
\end{enumerate}

\subsection{Návrat z prerušenia}
Ukončenie obsluhy prerušenia obnoví pôvodný vykonávaný kontext zo zásobníka:
\begin{enumerate}
  \item Zo zásobníka sa načíta a obnoví pôvodný stav registra príznakov:
  $$\reg{FREG} \leftarrow \text{Memory}[\reg{SREG}], \quad \reg{SREG} \leftarrow \reg{SREG} + 2$$
  (Obnovením registra \reg{FREG} sa zároveň bit \reg{FREG.INTR} vráti na \texttt{0}, čím sa uvoľní blokovanie, a hodnota \reg{FREG.MIE} sa vráti na pôvodný stav).
  \item Zo zásobníka sa obnoví čítač inštrukcií:
  $$\text{PC} \leftarrow \text{Memory}[\reg{SREG}], \quad \reg{SREG} \leftarrow \reg{SREG} + 2$$
  \item Procesor plynule pokračuje vo vykonávaní prerušeného toku inštrukcií v nasledujúcom takte bez akejkoľvek straty údajov.
\end{enumerate}

\section{Záver}
Architektúra \textbf{LIM M2} spája overenú stabilitu mikroprogramového riadenia s modernou efektivitou aktívnej dátovej cesty bez tretieho stavu. Odstránením zbernicových kolízií a parazitných kapacít jadro zaručuje kryštálovú čistotu digitálnych signálov a bezpečne odomyká frekvenčný rozsah od 200~MHz až po fyzické limity moderných kremíkových štruktúr.
\end{document}
